\begin{tikzpicture}[font=\small,
  brain/.pic={
    % brainstem and cerebellum first, so the cerebrum covers their tops
    \draw[thick, fill=white] (2.55,0.6) -- (2.45,-1.0) .. controls (2.6,-1.15) and (2.9,-1.15) .. (3.05,-1.0) -- (3.15,0.6);
    \draw[thick, fill=white] (1.45,0.3) ellipse (1.0 and 0.5);
    \draw[thick, fill=white] plot[smooth cycle, tension=0.7] coordinates
      {(0.1,1.9) (0.5,3.0) (1.5,3.8) (3.0,4.15) (4.5,3.9) (5.6,3.15) (6.05,2.1) (5.75,1.25)
       (5.0,0.95) (4.35,1.05) (4.15,0.6) (3.4,0.3) (2.2,0.45) (1.2,0.85) (0.35,1.3)};
    \fill[gray, fill opacity=0.10] plot[smooth cycle, tension=0.7] coordinates
      {(0.1,1.9) (0.5,3.0) (1.5,3.8) (3.0,4.15) (4.5,3.9) (5.6,3.15) (6.05,2.1) (5.75,1.25)
       (5.0,0.95) (4.35,1.05) (4.15,0.6) (3.4,0.3) (2.2,0.45) (1.2,0.85) (0.35,1.3)};
  },
  lab/.style={font=\small, align=center},
  lead/.style={thin}
]
% ================= outer (lateral) surface =================
\pic at (0,0) {brain};
\node[font=\small\bfseries] at (3.0,-1.75) {Outer (lateral) surface};
% lateral sulcus as a landmark
\draw[gray] (4.35,1.35) .. controls (3.6,1.6) and (3.0,1.9) .. (2.4,2.25);
% DLPFC and VLPFC
\fill[green!40] (4.95,3.05) ellipse[x radius=0.72, y radius=0.4, rotate=-38];
\fill[blue!35] (5.5,1.85) ellipse[x radius=0.52, y radius=0.3, rotate=-75];
\draw[lead] (5.2,3.2) -- (6.6,3.9) node[lab, right] {DLPFC};
\draw[lead] (5.65,1.75) -- (6.6,1.2) node[lab, right] {VLPFC};
\draw[-{Stealth[length=2mm]}, gray] (4.6,-1.0) -- (5.6,-1.0) node[right, font=\footnotesize, text=gray] {front};

% ================= inner (medial) surface =================
\begin{scope}[xshift=9.2cm]
\pic at (0,0) {brain};
\node[font=\small\bfseries] at (3.0,-1.75) {Inner (medial) surface};
% corpus callosum as a landmark
\draw[gray, line width=4pt, line cap=round]
  (1.75,2.2) .. controls (2.3,2.85) and (3.9,2.95) .. (4.4,2.55)
  .. controls (4.75,2.25) and (4.55,1.85) .. (4.15,1.95);
\node[font=\scriptsize, text=gray] at (2.85,2.1) {corpus callosum};
% dorsal ACC: band over the front of the corpus callosum
\draw[orange, line width=5pt, line cap=round]
  (2.9,3.15) .. controls (3.7,3.3) and (4.55,3.05) .. (4.8,2.55);
% ventral ACC: below and in front of the bend
\draw[purple!70, line width=5pt, line cap=round]
  (4.95,2.3) .. controls (5.0,1.95) and (4.85,1.65) .. (4.5,1.5);
% DMPFC
\fill[yellow!60] (4.75,3.55) ellipse[x radius=0.62, y radius=0.24, rotate=-28];
% amygdala, deep in the temporal lobe
\fill[red] (3.75,1.0) circle (0.13);
\draw[lead] (3.2,3.25) -- (2.3,4.65) node[lab, above] {Dorsal ACC};
\draw[lead] (4.95,3.6) -- (5.6,4.65) node[lab, above] {DMPFC};
\draw[lead] (4.95,1.95) -- (6.6,1.6) node[lab, right] {Ventral ACC};
\draw[lead] (3.8,0.88) -- (4.15,-0.3) node[lab, right] {Amygdala};
\draw[-{Stealth[length=2mm]}, gray] (4.6,-1.0) -- (5.6,-1.0) node[right, font=\footnotesize, text=gray] {front};
\end{scope}
\end{tikzpicture}
